
\subsection{Administración de Configuración y Documentación - Conservar versiones coherentes de nuestros artefactos [Pendientes: 1. Definir cambios, ramas, versiones, integración, nomenclatura y almacenamiento pesado. 2. Asignar administración, revisión, aprobación y restauración. 3. Comprobar servicio y recuperación y definir respaldo del repositorio de la aplicación.]}\label{sec:spmp-12-3}

La administración de configuración y documentación tiene como propósito conservar versiones identificables y coherentes de los productos desarrollados, facilitando su revisión, recuperación y relación con los requisitos y actividades del proyecto. Los elementos de configuración comprenden código, documentos, modelos, datos preparados y catálogo~\cite{adaptacion2026,std12207}.

Este proceso será aplicable a los entregables E1 a E7 descritos en la sección \ref{sec:spmp-6-4}. Su alcance incluye tanto los archivos utilizados para desarrollar la aplicación como aquellos necesarios para preparar, ejecutar y comprender las simulaciones FDS.

Para los documentos de Aldine, la documentación del servidor describe repositorios Git por proyecto y sincronización unidireccional hacia un repositorio consolidado en GitHub. El esquema conserva las fuentes y el PDF principal compilado y excluye repositorios internos y auxiliares de compilación; dispone de activación manual y de un temporizador aproximadamente cada cuatro horas. La identidad de los commits automáticos es Aldine y se distingue de la autoría individual de los contenidos.

\textbf{Identificación y evolución de los elementos de configuración.}

Cada artefacto deberá poder relacionarse con el entregable al que pertenece y con las actividades que originaron su creación o modificación. Para ello, se conservará información sobre su versión, procedencia, finalidad y relaciones con otros productos.

La siguiente tabla presenta los principales artefactos que estarán sujetos al proceso de administración de configuración, así como los momentos en los cuales deberán crearse o actualizarse.

\begin{longtable}{|p{0.23\dimexpr\linewidth-34pt\relax}|p{0.43\dimexpr\linewidth-34pt\relax}|p{0.34\dimexpr\linewidth-34pt\relax}|}
\caption{Evolución de los artefactos de documentación y código. Fuente: adaptación de estándares del proyecto~\cite{adaptacion2026}.}\\
\hline
\textbf{Artefacto} & \textbf{Creación o refinamiento} & \textbf{Información para mantener el artefacto actualizado y coherente}\\\hline
\endfirsthead
\hline
\textbf{Artefacto} & \textbf{Creación o refinamiento} & \textbf{Información para mantener el artefacto actualizado y coherente}\\\hline
\endhead
Requisitos (ISO/IEC/IEEE 29148) & Al precisar requisitos y al aprobar cambios durante el desarrollo. & Origen, versión y relaciones con diseño y prueba.\\\hline
Procesos del software (ISO/IEC/IEEE 12207) & Al planificar actividades y al registrar su ejecución. & Responsable, producto y evidencia.\\\hline
Modelado UML & En decisiones de arquitectura y cambios de interfaces o implementación. & Revisión de coherencia entre vistas.\\\hline
Pruebas (ISO/IEC/IEEE 29119) & Al planificar pruebas y después de cada ejecución o corrección. & Caso, entradas, versión y resultado.\\\hline
Calidad del software (IEEE 730) & Al preparar revisiones y tratar hallazgos de cada versión. & Criterio, problema, acción y reevaluación.\\\hline
Código de E5 & En cada incremento que modifique el producto. & Git y comprobaciones asociadas.\\\hline
Datos, entradas y resultados de E3-E4 & Al preparar, ejecutar, revisar o incorporar un escenario. & Procedencia, unidades, configuración e inventario.\\\hline
\end{longtable}

\textbf{Procedimiento de administración de configuración.}

Se propone utilizar el siguiente procedimiento para mantener la consistencia entre los artefactos:

\begin{enumerate}
    \item \textbf{Identificar los elementos.} Registrar los archivos y productos que forman parte de cada entregable, distinguiendo los documentos, el código fuente, los modelos, los datos de entrada y los resultados de simulación.

    \item \textbf{Controlar las modificaciones.} Documentar los cambios realizados y conservar su relación con los requisitos, las decisiones de diseño y las pruebas. Los cambios que afecten compromisos del proyecto deberán seguir los mecanismos de aprobación descritos en las secciones 6.6 y 10.1.

    \item \textbf{Revisar e integrar.} Comprobar que las modificaciones no generen inconsistencias con los demás componentes. La integración del código y de los documentos deberá considerar las comprobaciones correspondientes antes de establecer una nueva versión de referencia.

    \item \textbf{Identificar versiones.} Asignar versiones reconocibles a los productos y conservar la información que permita determinar qué archivos, requisitos y comprobaciones pertenecen a cada versión.

    \item \textbf{Conservar y respaldar.} Mantener los artefactos en los repositorios y espacios de almacenamiento acordados, procurando que exista una copia recuperable de la información necesaria para continuar el proyecto.

    \item \textbf{Recuperar y comprobar.} Cuando se identifique una pérdida, corrupción o modificación incorrecta, localizar una versión anterior, restaurar los archivos afectados y comprobar su integridad antes de continuar el trabajo.
\end{enumerate}

\textbf{Control de versiones e integración.}

Git será la herramienta utilizada para conservar la evolución del código fuente y de los artefactos que puedan administrarse mediante repositorios. Para el código de la aplicación se propone mantener una rama principal que represente una versión integrada y utilizar ramas de trabajo para desarrollar funcionalidades, resolver errores o modificar documentación.

Como convención inicial, podrían emplearse nombres descriptivos para las ramas, tales como \texttt{feature/}, \texttt{fix/} y \texttt{docs/}, acompañados de una identificación breve del cambio. Esta nomenclatura constituye una propuesta pendiente de aprobación por el equipo y no modifica el mecanismo de sincronización ya documentado para Aldine.

Antes de integrar cambios en la rama principal, se deberá revisar su correspondencia con los requisitos y comprobar que no introduzcan errores identificables. Los cambios deberán contar con una descripción que permita reconocer su finalidad, los archivos afectados y las verificaciones realizadas.

Se propone distinguir las versiones de trabajo de las versiones establecidas como referencia para revisiones, demostraciones y entregas. Las versiones finales de los entregables deberán identificarse de manera coherente, evitando que se utilicen simultáneamente archivos con modificaciones incompatibles.

\textbf{Administración de los datos y resultados de simulación.}

Los archivos de entrada y los resultados de FDS requieren especial atención debido a su tamaño y a la necesidad de reproducir las condiciones de ejecución.

Para cada escenario incorporado a la biblioteca se deberá conservar, como mínimo, la identificación del escenario, su ubicación de referencia, las condiciones ambientales utilizadas, la configuración de FDS, la versión del simulador y la ubicación de los resultados correspondientes. Cuando resulte pertinente, también se registrarán las transformaciones de coordenadas y las fuentes de datos empleadas.

Se propone almacenar los resultados de simulación de gran tamaño en un espacio separado del repositorio principal de código, conservando en este último o en el catálogo las referencias necesarias para localizarlos. El medio definitivo de almacenamiento deberá establecerse de acuerdo con el tamaño de la biblioteca, los recursos disponibles y las condiciones de acceso.

También se propone conservar un inventario de archivos y mecanismos de verificación de integridad, como sumas de comprobación, para detectar pérdidas o modificaciones accidentales. La organización del almacenamiento deberá permitir distinguir los archivos originales, los datos procesados, las entradas FDS y los resultados obtenidos.

\textbf{Respaldo y recuperación.}

La sincronización del servidor Aldine hacia GitHub contribuye a conservar versiones documentales, pero no reemplaza por sí sola la necesidad de comprobar la recuperación del servicio ni garantiza el respaldo de todos los archivos de la aplicación y de las simulaciones.

Por esta razón, se propone establecer copias de respaldo independientes de los repositorios y de los archivos necesarios para reconstruir el entorno de trabajo. El procedimiento deberá contemplar la identificación de la última versión disponible, la restauración de los archivos y la comprobación de su integridad.

Para el servidor documental se deberá verificar la posibilidad de recuperar los proyectos, las fuentes LaTeX y los archivos asociados. Para la aplicación, se comprobará la recuperación del código fuente y de las dependencias documentadas. En el caso de la biblioteca, se deberá verificar que los archivos restaurados correspondan con el inventario y el catálogo de escenarios.

Las frecuencias de respaldo, los destinos de almacenamiento y las pruebas periódicas de restauración deberán acordarse de acuerdo con la infraestructura disponible.

\textbf{Responsabilidades de administración y revisión.}

Daniel Ramírez es responsable de la administración y el mantenimiento del servidor documental Aldine, según la organización descrita en la sección 8.4.2.

La administración del repositorio de la aplicación, la revisión de modificaciones, la aprobación de integraciones y la ejecución de procedimientos de recuperación requerirán una asignación específica entre los integrantes del equipo. Estas responsabilidades deberán definirse considerando la organización del proyecto y los acuerdos de control de cambios de la sección \ref{sec:spmp-6-6}.

Hasta que se acuerden las convenciones definitivas de ramas, versiones, almacenamiento y respaldo, los procedimientos descritos constituyen una propuesta para su aplicación durante la Construcción y Semi-transición.
